% Appendix (IEEE: after the conclusion, before the acknowledgment).
\appendices
\section{First-Generation (v1) Training Images Rendered from LIDC-IDRI}
\label{app:v1}
The first (v1) suppressors, the lesion library and the synthetic test CTs come from LIDC-IDRI. This appendix describes how a CT is split
into exact components and projected, how a training sample is composed, and how lesions were used in v1.

\emph{CT decomposition.} Each CT was split into bone, non-lung soft tissue and lung. The body is the largest component of $\mathrm{HU}>-300$
(slice-wise hole filling, then a $1{\times}9{\times}9$ opening to detach the couch). Bone follows the iterative-relaxation recipe of the
reference method \cite[Algorithm~1]{angaitkar2026anatomy}, with a 10~mm support cube, the first pass starting from the reference mask and a
single closing after the ladder: a reference mask at 300~HU, thresholds $\tau\in\{200,150,100\}$~HU admitted only where at least 5\,\% of a
10~mm neighbourhood was already bone, union with the reference mask, a radius-2 closing to fill marrow, and removal of components smaller
than 1\,\% of the largest. Lung is air ($<-320$~HU) inside the body in components at least 20\,\% the size of the largest, closed and
hole-filled slice-wise (so vessels and nodules belong to the lung component) and without bone. Bone voxels were replaced in the
soft-tissue component by a normalised convolution of the surrounding soft tissue (Gaussian, 4~mm, with 12 and 36~mm fall-backs). Density is
water-equivalent, $\rho=\mathrm{clip}(1+\mathrm{HU}/1000,0,3)$~g\,cm$^{-3}$.

\emph{Polychromatic projection.} Each component was integrated along the anteroposterior axis (parallel beam), giving an areal density
$a_c$ (g\,cm$^{-2}$). The frame is a body crop: width 1.08 times the body width, height the lung extent plus 60~mm clipped to 0.75--1.2 of
the width (short scans cannot reach 1.2), area-resized to $512\times512$ and later squashed to a square exactly as real radiographs are.
Padded rows are set to air. The AP view is the horizontal mirror of the PA view. For a component $c$ with areal density $a_c$, a tube
spectrum $\phi(E)$ and an energy-integrating detector, a pixel value is
\begin{equation}
P_c=-\ln\frac{\sum_E \phi(E)\,E\,\exp\!\big(-(\mu/\rho)_c(E)\,a_c\big)}{\sum_E\phi(E)\,E},
\label{eq:physics}
\end{equation}
with ICRU-44 \cite{icru1989report44} cortical bone, soft tissue and lung attenuation (Table~4 of NIST SRD 126 \cite{hubbell2004xraymass})
and W-anode spectra (SpekPy \cite{poludniowski2021spekpy}, 12$^\circ$, 2.5~mm Al) at 70--140~kVp (training draws 80--130~kVp, the
synthetic evaluation uses 100~kVp). Components are stored as areal densities and the spectrum is applied at training time, so one render
serves every tube voltage. The sum of separately projected components differs from the projection of the whole body because of beam
hardening; we report this additivity gap (6.1\,\% at 100~kVp for 20~g\,cm$^{-2}$ of soft tissue and 3~g\,cm$^{-2}$ of bone).
Fig.~\ref{fig:pipeline} shows the decomposition and projection of one CT.

\begin{figure*}[!t]
\centering
\safefig[0.7\textwidth]{figures/fig_preprocessing_pipeline.pdf}
\caption{CT decomposition and polychromatic projection for one LIDC-IDRI test CT. (a) CT slice; (b) bone, lung and non-lung soft-tissue
masks; (c)--(e) areal densities of bone, non-lung soft tissue and lung; (f) composite radiograph $X$ at 100~kVp; (g) exact bone image $B$;
(h) exact soft-tissue image $S=X-B$; (i) a real radiograph (SZCH-X-Rays) for comparison.}
\label{fig:pipeline}
\end{figure*}

\emph{Training samples.} A training sample is composed on the GPU from the stored areal densities in this order: lesion insertion, one
random affine transform shared by all maps (rotation $\pm15^\circ$, scale 0.9--1.15, shift $\pm6\,\%$, mirror), the physics of
Eq.~\eqref{eq:physics}, component weighting and a monotone tone curve $\gamma$:
\begin{equation}
\begin{gathered}
S=\mathrm{norm}\big(P_{\mathrm{nl}}+g\,P_{\mathrm{lung}}\big),\quad B=\frac{w_b\,P_{\mathrm{bone}}}{p_{99}(P_{\mathrm{bone}})},\\
X=\frac{S+B}{1.5},\quad X'=\gamma(X),\quad S'=\gamma(S),\quad B'=X'-S',
\end{gathered}
\label{eq:compose}
\end{equation}
with $g\sim U(1,3)$ and $w_b\sim U(0.2,0.5)$. The tone curve is either a random gamma or histogram matching to the intensity quantiles of
a randomly drawn real radiograph (4\,000 NIH ChestX-ray14 training-list images). Defining $B'=X'-S'$ keeps the decomposition exact under
any monotone mapping. Gaussian blur is linear and applied identically to $X'$ and $B'$; quantum noise and scatter modify only $X'$. The
region of interest of the suppression loss (Eq.~\eqref{eq:loss}) comes from the lung component (top structures:
$P_{\mathrm{lung}}/p_{99.9}>0.3$). Unlike the Qure.ai projections, where lesions are inserted after the affine transform, lesions enter the
LIDC renders before the affine transform and the physics.

\emph{Lesion-faithful supervision in v1.} Each lesion (LIDC consensus of at least two readers, 50\,\% voxel agreement) was projected as its
excess density over the surrounding lung. The \emph{lesion-weighted loss} adds $(\mathrm{L1}+\mathrm{MSE})$ on a dilated lesion mask with
weight 1. \emph{Counterfactual lesion insertion (CLIA)} adds, before the physics, 0--3 lesion patches drawn from \emph{training} CTs at
random lung positions (including over bone) with random scale (0.7--1.5) and an amplitude of 0.6--1.6 times that scale; the bone target is
untouched, so the network is explicitly taught that a lesion on a rib is not bone. The v1 hypothesis H1 compares these arms on synthetic
LIDC test CTs with exact truth (Section~\ref{sec:res-fidelity}).
